\documentclass[11pt]{amsart}
\usepackage[T1]{fontenc}
\usepackage[utf8]{inputenc}
\usepackage{amssymb}
\usepackage[margin=1.05in]{geometry}
\usepackage[hidelinks]{hyperref}
\newcommand{\stamptime}{\the\numexpr\time/60\relax:\ifnum\numexpr\time-60*(\time/60)\relax<10 0\fi\the\numexpr\time-60*(\time/60)\relax}
\AddToHook{shipout/foreground}{\put(0.5\paperwidth,-\paperheight+0.55cm){\makebox[0pt]{\tiny\ttfamily unit/doc/\jobname.tex \ --- \ compiled \today, \stamptime}}}
\newtheorem{theorem}{Theorem}
\newtheorem{proposition}[theorem]{Proposition}
\newtheorem{lemma}[theorem]{Lemma}
\newtheorem{corollary}[theorem]{Corollary}
\theoremstyle{definition}
\newtheorem{definition}[theorem]{Definition}
\theoremstyle{remark}
\newtheorem{remark}[theorem]{Remark}
\newcommand{\Om}{\Omega}
\newcommand{\Map}{\mathrm{Map}_*}
\newcommand{\Z}{\mathbb{Z}}
\newcommand{\code}[1]{\texttt{#1}}
\newcommand{\cT}{\mathsf{isTrunc}}
\newcommand{\cC}{\mathsf{isConn}}
\newcommand{\refl}{\mathsf{refl}}
\newcommand{\ap}{\mathsf{ap}}
\newcommand{\el}{\mathsf{el}}
\newcommand{\hc}[2]{[#1,#2]_*}
\newcommand{\Fib}{\mathrm{Fib}}
\newcommand{\fib}{\mathrm{fib}}
\newcommand{\tr}[1]{\|#1\|}
\newcommand{\cls}[1]{[#1]}
\title{A Postnikov presentation of $\|S^2\|_3$, and $\pi_3(S^2)=\Z$}
\date{28 October 2026}
\begin{document}
\maketitle

\noindent\textbf{Purpose.} Informal proof for \S2 of
\code{doc/pi3-s2-postnikov-handoff-8.txt} (\S3d of the plan). Let
$T=\tr{S^2}_3$ be the constructed 3-truncation
(\code{SphereTruncation.SphereTrunc.trunc3}). We prove
\[ \tr{\textstyle\sum_{B:\mathrm{Ab}}\sum_{k:K(\Z,2)\to_*K(B,4)}\ \Fib k\simeq_*T}\]
(Theorem~\ref{t:pres}); with \code{MapFib.Final.rep-pteq} it gives the
mapping theorem at $X=T$ (Corollary~\ref{c:mapT}), and a two-instance form
of \code{SphereUniversal.Assemble} then gives $\pi_3(S^2)\cong\Z$ from the
sphere axiom alone (Theorem~\ref{t:final}). Written before the formal
proof; numbered statements are meant to be declarations. Foundations as in
the development (univalence, images; function extensionality derived); no
new axiom. The group $B$ is \emph{not} assumed to be $\Z$: it is
$\Om^3$ of the fibre over the base point, and $B\cong\Z$ is a consequence of
the final comparison.

\smallskip\noindent\textbf{The idea.} (1) The constructed 2-truncation
$V=\tr T_2$ is pointed equivalent to $K(\Z,2)$: both have the universal
property of the 2-truncation of $S^2$ for maps into 2-types
(\code{Trunc3.up}, \code{TruncationProps.PtUP.up}, \code{MapFib.SourceS.Sig.sigma-star}).
(2) The fibre $F_0$ of the unit $u:T\to V$ over the base point is a
2-connected 3-type, hence $K(B,3)$ with $B=\Om^3F_0$ (a recognition
principle, from the delooping corollaries of \code{DeloopMaps} and
Whitehead's principle). (3) In the library $K(B,4)$ \emph{is} the type of
$K(B,3)$-torsors (\code{EM}: $K(A,n{+}3)=T\,K(A,n{+}2)$, \code{Deloop}),
i.e.\ of types $Y$ with $\tr Y$ and a homogeneous structure
$\mu:(y:Y)\to K(B,3)\simeq_*(Y,y)$; and the fundamental theorem of
\code{Deloop} identifies $(\mathrm{pt}=t)$ with the underlying type of $t$.
So a map $k:V\to K(B,4)$ with fibres the fibres of $u$ is the same as a
homogeneous structure on every fibre of $u$. (4) The homogeneous
structures on the fibres of $u$ form a family of \emph{sets} over the
simply connected total space $T$; so the one structure given at the base
point by (2) extends to all of $T$. This is where simple connectivity
removes the twisting: the choices at a point form a set (an
$\mathrm{Aut}(B)$-torsor), and a family of sets over a 1-connected type has
a section through any given point.

\section{Conventions}

As in \code{doc/mapping-fib.pdf}: \code{isTrunc n X} is ``$X$ is
$(n-2)$-truncated'', \code{isConn n X} is ``$X$ is $(n-2)$-connected'';
$K^A_m=K(A,m)$ is \code{EMMaps.KK[A].K m}, pointed by \code{pt m}, with
$e_m:\Om K^A_{m+1}\simeq_*K^A_m$ (\code{KK.e}; $e_0$ is the band of
$K(A,1)$, a group isomorphism); $K_2=K^\Z_2$ (\code{KZ2}); $F;G$ is
\code{ptcomp F G}, $(F;G)_1=G_1\circ F_1$ by definition; $\fib_f(z)$ is
\code{Whitehead.fib}; $\Fib k$ is \code{Fibres.Fib}. $(T,t_0,\eta)$ is
\code{trunc3}: $T$ is \code{isTrunc 5} and \code{isConn 3}, and
$\eta^*:\Map(T,Y)\to\Map(S^2,Y)$ is an equivalence for $Y$ \code{isTrunc 5}
(\code{up}). $\sigma:S^2\to_*K_2$ is \code{MapFib.SourceS.sigma}, and
$\sigma^*$ is an equivalence on maps into every \code{isTrunc 4} type
(\code{Sig.sigma-star}). $V=\code{(TO three).T}\,T$ is \code{isTrunc 4},
$u=\code{PtUP[four, TO three, T, t0].eta}:T\to_*V$, $v_0=u\,t_0$, the
fibres of $u$ are \code{isConn 4} (\code{TruncOp.conn}), and $u^*$ is an
equivalence on maps into \code{isTrunc 4} types (\code{PtUP.up}).

\section{$V$ is $K(\Z,2)$}

\begin{definition}\label{d:maps}
$r=(\eta^*)^{-1}\sigma:T\to_*K_2$ (\code{up} at $Y=K_2$, which is
\code{isTrunc 4}, hence 5), so $\eta;r=\sigma$;
$\varphi=(u^*)^{-1}r:V\to_*K_2$, so $u;\varphi=r$;
$\psi=(\sigma^*)^{-1}(\eta;u):K_2\to_*V$, so $\sigma;\psi=\eta;u$.
\end{definition}

\begin{lemma}\label{l:inv}
(a) $\psi;\varphi=\mathrm{id}_{K_2}$; (b) $r;\psi=u$; (c)
$\varphi;\psi=\mathrm{id}_V$. All equalities of pointed maps.
\end{lemma}
\begin{proof}
Each by injectivity of an equivalence (\code{Equivalences.injective}) and
associativity (\code{ptcomp-assoc}).
(a) $\sigma^*$ at $Y=K_2$: $\sigma;(\psi;\varphi)=(\sigma;\psi);\varphi
=(\eta;u);\varphi=\eta;(u;\varphi)=\eta;r=\sigma=\sigma;\mathrm{id}$
(\code{ptcomp-idr}).
(b) $\eta^*$ at $Y=V$ (\code{isTrunc 4}, hence 5):
$\eta;(r;\psi)=(\eta;r);\psi=\sigma;\psi=\eta;u$.
(c) $u^*$ at $Y=V$: $u;(\varphi;\psi)=(u;\varphi);\psi=r;\psi=u=u;\mathrm{id}$, by (b).
\end{proof}

\begin{proposition}\label{p:psi}
$\psi_1:K_2\to V$ is an equivalence, with inverse $\varphi_1$.
\end{proposition}
\begin{proof}
The first components of Lemma~\ref{l:inv}(a),(c) are
$\varphi_1\circ\psi_1=\mathrm{id}$ and $\psi_1\circ\varphi_1=\mathrm{id}$
(by definition of \code{ptcomp}); \code{happly} and \code{qinv}.
\end{proof}

\begin{remark}
We never use the fibres of $r$: the presentation is built over $V$, where
the fibres of $u$ are 2-connected by construction of the truncation, and
transported to $K_2$ along $\psi$ (\S\ref{s:pres}). (That the fibres of
$r$ are 2-connected follows, since $r=\varphi\circ u$ with $\varphi$ an
equivalence, but it is not needed.)
\end{remark}

\section{Recognising $K(B,3)$}

\begin{lemma}[delooping an equivalence]\label{l:deloop}
Let $(A,a_0)$, $(Y,y_0)$ be 0-connected, $f:\Om A\to_*\Om Y$ with $f_1$ an
equivalence, and suppose $\mathrm{Unloop}\,f$ is contractible, with centre
$(G,\gamma:f=\Om G)$. Then $G$ is a pointed equivalence $A\simeq_*Y$.
\end{lemma}
\begin{proof}
$(\Om G)_1$ is pointwise equal to $f_1$ (\code{happly} of
$\ap\,(\lambda w.\,w_1)\,\gamma$), so it is an equivalence (\code{E.homotopy}, as
in \code{EM1.mu-loop-equiv}); Whitehead's principle
(\code{Whitehead.Pointed.equiv}) concludes. This is the argument of
\code{EM.EM1.mu}, stated once.
\end{proof}

\begin{definition}[$\Om^3$ as an abelian group]\label{d:B}
For a pointed 3-type $(F,f_0)$ let $B_F$ be the abelian group with carrier
$\Om^3F=(\refl_{\refl f_0}=_{\Om^2F}\refl_{\refl f_0})$, a set since $F$ is
\code{isTrunc 5}; zero $\refl$, addition path composition in $\Om^2F$,
inverse \code{sym}; associativity \code{tassoc}, right unit by definition
of \code{trans}, right inverse \code{trans-sym-r}, commutativity
Eckmann--Hilton (\code{EH[$\Om F$].eh}: 2-loops of $(\Om F,\refl)$ commute).
\end{definition}

\begin{proposition}[recognition]\label{p:rec}
Let $(F,f_0)$ be \code{isConn 4} and \code{isTrunc 5}, $B=B_F$. Then there
is a pointed equivalence $e_F:K^B_3\simeq_*(F,f_0)$.
\end{proposition}
\begin{proof}
Three applications of Lemma~\ref{l:deloop}, climbing
$\Om^2F\to\Om F\to F$. Levels: $\Om^2F$ is \code{isConn 2}, \code{isTrunc
3}; $\Om F$ is \code{isConn 3}, \code{isTrunc 4}; $F$ is \code{isConn 4},
\code{isTrunc 5} (\code{conn-id}, \code{conn-down}; identity types of a
truncated type).
\emph{Step 1} ($G_1:K^B_1\simeq_*\Om^2F$). $f_1=(e_0)_1:\Om K^B_1\to_*B=\Om(\Om^2F,\refl)$,
the band, pointed by \code{band-unit}; it is multiplicative
(\code{band-hom}: addition of $B$ is composition in $\Om^2F$ by
Definition~\ref{d:B}), so $\mathrm{Unloop}\,f_1$ is contractible by
Corollary~12 (\code{Mult.cor12}: $K^B_1$ 0-connected, $\Om^2F$ 1-truncated).
\emph{Step 2} ($G_2:K^B_2\simeq_*\Om F$). $f_2=(e_1)_1;G_1:\Om K^B_2\to_*\Om^2F=\Om(\Om F,\refl)$,
an equivalence (composite); $\mathrm{Unloop}\,f_2$ contractible by
Corollary~11 with $c=1$ ($K^B_2$ is \code{isConn 3}, $\Om F$ is
\code{isTrunc 4} $=$ \code{isTrunc (2+(2+0))}).
\emph{Step 3} ($G_3:K^B_3\simeq_*F$). $f_3=(e_2)_1;G_2$; Corollary~11 with
$c=2$ ($K^B_3$ is \code{isConn 4}, $F$ is \code{isTrunc 5}, hence
\code{isTrunc 6} $=$ \code{isTrunc (3+(3+0))}).
In each step both sides are 0-connected.
\end{proof}

\section{Homogeneity of the fibres of $u$}

Write $E_z=\fib_u(z)$ for $z:V$; $F_0=E_{v_0}$ pointed by
$f_0=(t_0,\refl)$. Each $E_z$ is \code{isConn 4} (\code{TruncOp.conn}) and
\code{isTrunc 5} ($\Sigma$ of $T$ and of identity types of the 2-type $V$).
Let $B=B_{F_0}$ (Definition~\ref{d:B}), $K_3=K^B_3$, $k_3$ its point, and
$e_0=e_{F_0}:K_3\simeq_*(F_0,f_0)$ (Proposition~\ref{p:rec}).

\begin{lemma}[sections of a family of sets over a 1-connected type]\label{l:ext}
Let $(X,x_0)$ be \code{isConn 3}, $Q:X\to U$ a family of sets and $q_0:Q\,x_0$.
Then there is $s:(x:X)\to Q\,x$ with $s\,x_0=q_0$.
\end{lemma}
\begin{proof}
Lemma~4 (\code{lemma4} with $t=0$, $c=2$: $X$ is \code{isConn 3}, each
$Q\,x$ is \code{isTrunc (2+0)}): the pointed sections
$\mathrm{PtSec}(X,x_0,Q,q_0)$ form a contractible type; $s$ and $s\,x_0=q_0$
are the components of its centre. (Suggested by the referee, in place of
a unique-choice argument.)
\end{proof}

\begin{lemma}\label{l:setfam}
For every $t:T$, $Q\,t=\bigl(K_3\simeq_*(E_{u t},(t,\refl))\bigr)$ is a set.
\end{lemma}
\begin{proof}
Corollary~5(ii) (\code{cor5-equiv} with $t=1$, $c=3$): $K_3$ is
\code{isConn 4} and $E_{ut}$ is \code{isTrunc (3+2)}.
\end{proof}

\begin{definition}[the homogeneous structure]\label{d:mu}
$\mu':(t:T)\to Q\,t$ is the section of Lemma~\ref{l:ext} for $X=T$ (which is
\code{isConn 3}), $Q$ of Lemma~\ref{l:setfam}, and $q_0=e_0$ (note
$E_{u t_0}=F_0$ and $(t_0,\refl)=f_0$ by definition). For $z:V$ and
$y=(t,p):E_z$, $\mu_z(y):K_3\simeq_*(E_z,y)$ is $\mu'(t)$ moved along $p$
($J$ on $p:u\,t=z$). So $\mu_z:\mathrm{Mu}(E_z)$ in the sense of
\code{Deloop.Torsors[$K_3$,$k_3$]}.
\end{definition}

\section{The classifying map}

Let $\mathcal{T}=\code{Deloop.Torsors[$K_3$,$k_3$]}$: $\mathrm{Tor}=\sum_{k:\mathrm{BAut}}\tr{\el k}\times\mathrm{Mu}(\el k)$,
where $\mathrm{BAut}$ is the small image of the point $K_3$ in $U$ and
$\el:\mathrm{BAut}\to U$ its inclusion. By definition of \code{EM}
(\code{EMn.Kstage}, \code{Step.next}), $K^B_4$ \emph{is} $\mathrm{Tor}$, its
point is the centre $\mathrm{pt}$ of \code{Step[1].D}, and the fundamental
theorem \code{D.F} gives equivalences
$\tau_t:(\mathrm{pt}=t)\simeq\el(t_1)$ (\code{F.tr}, \code{F.equiv}).
These by-definition facts are stated once, by \code{refl}.

\begin{lemma}[torsor from a homogeneous structure]\label{l:tor}
Let $Y:U$ with $\tr Y$ and $\mu:\mathrm{Mu}(Y)$. There is a torsor
$\langle Y,\mu\rangle:\mathrm{Tor}$ and a path $\pi_Y:\el(\langle Y,\mu\rangle_1)=Y$.
\end{lemma}
\begin{proof}
$\mathcal{T}.\code{Fib}\,Y=\sum_{k:\mathrm{BAut}}(\el k=Y)$ is a proposition
(\code{fib-prop}), inhabited at every point $y$ of $Y$
(\code{fib-contr Y ($|y|$, $\mu$, $y$)}), hence inhabited, by induction on
$\tr Y$ (at the universe level of $\mathcal{T}.\code{Fib}\,Y$). Let $(k,\pi_Y)$ be
this element; $\langle Y,\mu\rangle=(k,\ \pi_Y^{-1}{}_*(\tr Y\text{-witness},\mu))$, as in
\code{homog}.
\end{proof}

\begin{definition}[$k'$]\label{d:kp}
$k'_1:V\to K^B_4$, $k'_1(z)=\langle E_z,\mu_z\rangle$ (with $\tr{E_z}$ from
\code{conn-inh}). Let $c_z=\mathrm{idtoeqv}(\pi_{E_z}):\el(k'_1(z)_1)\simeq E_z$.
The pointing is $k'_0=\bigl(\tau^{-1}_{k'_1 v_0}(c_{v_0}^{-1}f_0)\bigr)^{-1}:k'_1v_0=\mathrm{pt}$.
\end{definition}

\begin{proposition}\label{p:tot}
$\Fib k'=\sum_{z:V}(k'_1z=\mathrm{pt})\simeq T$.
\end{proposition}
\begin{proof}
A composite of four equivalences:
$\sum_z(k'_1z=\mathrm{pt})\simeq\sum_z(\mathrm{pt}=k'_1z)$ (\code{sym}, fibrewise),
$\simeq\sum_z\el(k'_1(z)_1)$ ($\tau$, fibrewise),
$\simeq\sum_zE_z$ ($c_z$, fibrewise; \code{Fibres.sigma-equiv} for the three),
$\simeq T$ ($(z,(t,p))\mapsto t$, inverse $t\mapsto(u\,t,(t,\refl))$; $J$ on $p$).
\end{proof}

The pointing of $k'$ plays no role in Proposition~\ref{p:tot}: the
underlying type of $\Fib k'$ does not depend on it.

\section{The presentation}\label{s:pres}

\begin{definition}
$k_T=\psi;k':K_2\to_*K^B_4$, so $(k_T)_1=k'_1\circ\psi_1$ by definition.
\end{definition}

\begin{proposition}\label{p:fibkT}
$\Fib k_T\simeq T$.
\end{proposition}
\begin{proof}
$\Fib k_T=\sum_{x:K_2}(k'_1(\psi_1x)=\mathrm{pt})\simeq\sum_{z:V}(k'_1z=\mathrm{pt})$
by $\Sigma$ over the equivalence $\psi_1$ of bases
(\code{SphereAxiom.sigma-base}, Proposition~\ref{p:psi}); then
Proposition~\ref{p:tot}.
\end{proof}

\begin{theorem}[a Postnikov presentation of $T$]\label{t:pres}
For $B=B_{F_0}$, $k_T$ and the equivalence $e:\Fib k_T\simeq T$ of
Proposition~\ref{p:fibkT}: merely $e(\text{base point of }\Fib k_T)=t_0$, and
every such path makes $e$ a pointed equivalence $\Fib k_T\simeq_*(T,t_0)$.
In particular
$\tr{\sum_{B:\mathrm{Ab}}\sum_{k:K_2\to_*K^B_4}\Fib k\simeq_*(T,t_0)}$.
\end{theorem}
\begin{proof}
$T$ is 0-connected (\code{conn-path}). The truncated $\Sigma$ over
$\mathrm{Ab}$ lives in the next universe ($\mathrm{Ab}$ has a field in $U$);
the formal statement is the first form, which avoids it (referee), and the
mere path is eliminated where it is used, into a proposition
(Corollary~\ref{c:mapT}).
\end{proof}

Everything but the last pointing is a construction; nothing is chosen.

\section{The mapping theorem at $T$ and $\pi_3(S^2)$}

\begin{corollary}[\code{MapThm} at $X=T$]\label{c:mapT}
For $P=\Fib\kappa$ ($\kappa$ the cup square, \code{H4Proof.h4}) and
$s=\code{MapFib.Final.s}$, $s^*:\hc PT\to\hc{S^2}T$ is a bijection.
\end{corollary}
\begin{proof}
Being an equivalence is a proposition (\code{EquivProp.isEquiv-prop});
eliminate the truncation of Theorem~\ref{t:pres} and apply
\code{rep-pteq} with $(A,B,k)=(\Z,B,k_T)$.
\end{proof}

\begin{proposition}[Theorem~B, two instances]\label{p:two}
\code{SphereUniversal.Assemble} only uses \code{Rep.rep} at $X=P$
($s^*$ on $\hc PP$) and $X=T$ ($s^*$ on $\hc PT$). Restated with these two
bijections as hypotheses (a module \code{Two} taking $T$-data, $(P,p_0)$,
\code{isTrunc 5} $P$, $s$, \code{sP}, \code{sT}), the proof is the same
text; \code{Assemble} becomes its instance at
$\code{sP}=\code{rep}\,P$, $\code{sT}=\code{rep}\,T$.
\end{proposition}

\begin{theorem}[$\pi_3(S^2)=\Z$]\label{t:final}
For $(S^2,\mathsf{base},\mathsf{surf})$ with the sphere axiom,
$\pi_3(S^2)\cong\Z$ as groups.
\end{theorem}
\begin{proof}
Proposition~\ref{p:two} with $T$ from \code{trunc3}, $P=\Fib\kappa$,
$s$ as above, \code{sP} $=$ \code{rep-fib} for the triple $(\Z,\Z,\kappa)$,
\code{sT} from Corollary~\ref{c:mapT}, gives $\pi_3S^2\cong\pi_3P$; and
$\pi_3P\cong\Z$ (\code{PostnikovRealization.Triple.pi3-iso}).
\end{proof}

\begin{remark}[on $B$]
The proof gives a group $B$ with merely $T\simeq_*\Fib k_T$, and the
comparison then shows $\pi_3T\cong\pi_3P\cong\Z$; since $\pi_3\Fib k_T\cong B$
(\code{pi3-iso}), \emph{merely} $B\cong\Z$ follows, a posteriori (the
pointing of $e$ could be made explicit, which would remove ``merely''). It
is not used.
\end{remark}

\section{Correspondence with the formal proof}
\emph{Status: checked} (\code{./segc sag/pi3-s2-postnikov all}: 71 modules,
0 failed; a fresh process from the directory: 3135 declarations, 0 failed,
1.3\,s; every new declaration under 0.02\,s). Module \code{Presentation}
(entries in \code{present.opens}, 103, every line ``nothing else opened'';
checked against this note by kind, not entry by entry, see its header):
\code{deloop-equiv} (Lemma~\ref{l:deloop}); \code{O3[F,f0,tF]}: \code{X1},
\code{x1}, \code{A3}, \code{add}, \code{set}, \code{ab}
(Definition~\ref{d:B}); \code{Rec[F,f0,cF,tF]}: \code{cY2}, \code{cX1},
\code{tX1}, \code{f1}, \code{cell1}, \code{G1}, \code{f2}, \code{cell2},
\code{G2}, \code{f3}, \code{cell3}, \code{e3} (Proposition~\ref{p:rec};
step~1 is stated on \code{EM1.K1}, the definition of $K(B,1)$);
\code{Pres[S2,base,surf,ump,TR]} for any \code{Trunc3} \code{TR}: \code{V},
\code{v0}, \code{tV}, \code{tK2}, \code{r}, \code{r-eq}, \code{phi},
\code{phi-eq}, \code{psi}, \code{psi-eq} (Definition~\ref{d:maps}),
\code{psi-phi}, \code{r-psi}, \code{phi-psi} (Lemma~\ref{l:inv}),
\code{psi-equiv} (Proposition~\ref{p:psi}); \code{E}, \code{e-conn},
\code{e-trunc}, \code{f0}, \code{RF}, \code{Q}, \code{q-set}
(Lemma~\ref{l:setfam}), \code{ext} (Lemma~\ref{l:ext}), \code{mu'},
\code{mu} (Definition~\ref{d:mu}); \code{DS} (the instance
\code{EM.Step[1, EMn[B].Kstage 1]}, as the referee suggested),
\code{K4-def}, \code{pt4-def}, \code{K3-def} (by \code{refl}), \code{tau};
\code{tor-fib}, \code{tor}, \code{tor-el} (Lemma~\ref{l:tor}); \code{k1},
\code{c}, \code{kp} (Definition~\ref{d:kp}); \code{sym-equiv},
\code{tot-ret}, \code{tot-fib}, \code{tot} (Proposition~\ref{p:tot});
\code{kT}, \code{kT-def}; \code{fib-kT} (Proposition~\ref{p:fibkT});
\code{pres-pt}, \code{pres} (Theorem~\ref{t:pres});
\code{Final[S2,base,surf,ump]}: \code{sT} (Corollary~\ref{c:mapT}),
\code{A2}, \code{pi3-S2 : GIso[pi3 S2 base, ab-grp Zab]}
(Theorem~\ref{t:final}). Module \code{SphereUniversal}: \code{Two}
(Proposition~\ref{p:two}), \code{Assemble} now its instance
(\code{pi3.opens} updated).

\emph{Remaining extras} (\code{present.opens} header): \code{Rec.cell1}
opens the implementation of the band when \code{band-hom} is checked as
the multiplicativity of $f_1$ (a checker-defect candidate; the text uses
only that addition in $B$ is composition); the opaque \code{TO} is opened
to the type of the truncation in the declarations about the fibres $E_z$
(the etaV mechanism of \code{fibre.opens}); \code{Final.sT} opens the
carrier of classes (extra (ii) of \code{doc/pi3-s2-postnikov.pdf}).
\end{document}
